\section{Síntesis y ampliación}
Esta entrevista con Graham puede resumirse en tres temas que avanzan en paralelo: arquitectura, evaluación y gobernanza. Cada uno de estos caminos requiere una práctica de ingeniería intensiva, una evaluación integral y una reflexión sobre seguridad para que los agentes sean ampliamente adoptados en el desarrollo de software de alto riesgo.

\begin{table}[H]
\centering
\begin{tabular}{p{0.34\textwidth} p{0.58\textwidth}}
\toprule
Tema & Puntos clave \\
\midrule
Contexto y motivación & El software está devorando todas las industrias; delegar la capacidad de «escribir código» a los agentes permite que más equipos aborden los problemas con una mentalidad de ingeniería. \\
\midrule
Arquitectura y evaluación & OpenHands ofrece una plataforma de agentes transparente, y SWE-bench Lite/Full ofrece métricas de evaluación de ingeniería con dificultad progresiva. \\
\midrule
Futuro y gobernanza & Se necesitan datos con estilo de agente, más dimensiones de evaluación y mecanismos de seguridad integrados para que los agentes puedan incorporarse con seguridad al flujo de desarrollo convencional. \\
\bottomrule
\end{tabular}
\caption{Los tres elementos centrales de esta charla.}
\end{table}

Lecturas adicionales:
\begin{itemize}
    \item Repositorio de código abierto de OpenHands: \href{https://github.com/OpenHands/OpenHands}{https://github.com/OpenHands/OpenHands}.
    \item Artículo de SWE-bench: Jimenez et al., ``SWE-bench: Can Language Models Resolve Real-World GitHub Issues?'' ICLR 2024.
    \item Blog oficial de All Hands AI y su canal de retroalimentación de la comunidad, enfocados en la seguridad de los agentes y la colaboración humano-máquina.
\end{itemize}

\subsection{Resumen de la sección}
\begin{itemize}
    \item Las tres líneas que enfatiza esta charla —arquitectura, evaluación y gobernanza— deben avanzar en paralelo para llevar a los agentes de forma segura a producción.
    \item Seguir la evolución de OpenHands y SWE-bench ayuda a llevar el desarrollo de software agéntico a la práctica.
\end{itemize}

\end{document}
